\documentclass[12pt]{article}
\usepackage[margin=1in]{geometry}
\usepackage[T1]{fontenc}
\usepackage{helvet}
\renewcommand{\familydefault}{\sfdefault}
\usepackage{hyperref}
\usepackage{xcolor}
\usepackage{enumitem}
\usepackage{titlesec}

\definecolor{linkblue}{HTML}{00D8FF}
\newcommand{\hlink}[2]{\href{#1}{\textcolor{linkblue}{\underline{#2}}}}
\hypersetup{colorlinks=false}

\titleformat{\section}{\large\bfseries}{}{0em}{}
\titlespacing{\section}{0pt}{16pt}{8pt}

\setlist[itemize]{leftmargin=1.4em, itemsep=8pt, topsep=4pt}

\pagestyle{plain}

\begin{document}

\noindent\textbf{Shadrack Timothy John} is a \textbf{full-stack developer and AI/ML engineer} with over 4 years of experience building production software and AI-driven tools for clients across Tanzania. He works as an indie developer under \hlink{https://antera.co.tz}{ANTERA} (Advanced Neural Technologies and Engineering Research Agency), a practice delivering modern web systems, applied generative AI, and data infrastructure. His work bridges applied software engineering and independent AI research, with a focus on how low-resource languages such as Kiswahili are represented in modern language models. He is a fast learner, disciplined, and committed to professional growth, and he is looking for a role where that mix of engineering and research can be put to work.

\section*{Professional Profile}
\begin{itemize}
    \item \textbf{Academic Background:} He is a Computer Engineering Graduate as of 2026 from the \hlink{https://www.udom.ac.tz}{University of Dodoma}, following a prior Diploma in Computer Science from the University of Dar es Salaam.
    \item \textbf{Areas of Expertise:} His work spans \textbf{Full-Stack Web Development}, \textbf{Generative AI and LLMs}, \textbf{Data Engineering}, and \textbf{IoT and Edge Computing}.
    \item \textbf{Key Practice:} He delivers production software for clients across Tanzania as an indie developer under \hlink{https://antera.co.tz}{ANTERA}, serving SMEs, organizations, and independent research projects.
    \item \textbf{Digital Presence:} He publishes technical write-ups and project work at \hlink{https://antera.co.tz/blog}{antera.co.tz/blog} and \hlink{https://zuck30.github.io/zuck30}{zuck30.github.io/zuck30}.
\end{itemize}

\vspace{6pt}
\noindent Shadrack's work at ANTERA sits at the practical end of modern software delivery: production web engineering held together with a taste for clean user interfaces and thoughtful design. The result is software that ships, gets used, and still looks considered.

\section*{Skills and Methodology}
\noindent He builds and ships across the full stack, applying the same approach to client work and to his own research:

\begin{itemize}
    \item \textbf{Full-Stack Delivery:} He builds with \hlink{https://react.dev}{React.js}, Next.js, Tailwind CSS, and TypeScript on the frontend, and Django and FastAPI on the backend, deploying to Render, Netlify, Vercel, Supabase, Cloudflare, or Netpoa, or whichever platform gets the job done cleanly and without friction.
    \item \textbf{Applied AI:} A core part of his work is Python, backend prompt engineering, \textbf{fine-tuning} (LoRA, QLoRA), and \textbf{RAG systems}, used to build production chatbots and AI features backed by Hugging Face models and ChromaDB.
    \item \textbf{Data Engineering:} He designs data pipelines and cleaning workflows across PostgreSQL, MySQL, and SQLite, integrating APIs to automate reporting for clients.
    \item \textbf{Working Style:} He communicates with integrity, respect, and empathy, organizes people around a goal without bias, and takes a task from requirements through to deployment and maintenance.
\end{itemize}

\noindent Shadrack's background bridges \textbf{applied software engineering} and \textbf{independent AI research}, with a specific focus on how low-resource languages are represented in modern language models.

\section*{Experience and Research}
\begin{itemize}
    \item \textbf{Indie Developer, \hlink{https://antera.co.tz}{ANTERA} (2023 to Present):} His practice has delivered production web and AI systems for clients across Tanzania, including \hlink{https://swahiba.vercel.app}{Swahiba}, a Kiswahili-English chat frontend, taking each project from requirements through deployment and maintenance, alongside AI features backed by Hugging Face models, ChromaDB vector retrieval, and fine-tuned open-weight LLMs.
    \item \textbf{Swahili LLM Research:} He built a \textbf{70M-parameter Transformer} language model from scratch for Kiswahili-English mixed-language communication, training a custom SentencePiece tokenizer (1,400 vocab) on 200K synthetic samples ($\sim$30M tokens) and reducing training loss from 10.0 to 2.5 over 8 epochs on consumer hardware.
    \item \textbf{Malaika, AI Companion:} He fine-tuned Mistral-7B and Qwen2.5-VL-7B-Instruct via Hugging Face, integrated ChromaDB for long-term memory retrieval, and added Moondream2 for real-time visual awareness.
    \item \textbf{Frontend and Product Engineering Intern, \hlink{https://ictpack.com}{ICTPack Solutions Ltd}:} He ships ICTPack's client-facing platforms end to end, from auth and payment flows to pixel-level UI polish, across farm tracking, cybersecurity, and cross-border cargo products, usually as the sole or lead contributor on each repository. His work spans integration (JWT, OTP, and PIN auth flows, REST APIs, Google Maps, Selcom and CIPS wallets and payments, and Supabase edge functions behind the DeepSeek-powered ``Miss.\ ICTPack'' AI assistant), Swahili localization and Kiswahili defaults for the Tanzanian market, and delivery tooling (GitLab CI pipelines, README guides, and Docker configuration).
    \item \textbf{NOC Technical Intern, \hlink{https://azampay.com/}{Azampay Tanzania} (Jul 2023 to Sep 2023):} He took part in agile development cycles and code review, and worked to reduce response time for recurring technical faults. Alongside fellow developer Andrew Malisa, he developed an internal CLI tool that collects logs, statuses, and incidents of cron jobs.
\end{itemize}

\section*{Educational Qualification}
\begin{itemize}
    \item \textbf{2022 to 2026:} University of Dodoma, BSc in Computer Engineering.
    \item \textbf{2020 to 2022:} College of ICT, University of Dar es Salaam, Diploma in Computer Science.
\end{itemize}

\section*{Standing and Interests}
\begin{itemize}
    \item \textbf{Open To:} Full-time, remote, or contract roles in software engineering, AI/ML, or data engineering, and open-source collaboration on projects for East African languages, contexts, and problems.
    \item \textbf{Active Interests:} His recent focus areas include \textbf{Low-Resource Language Modeling}, \textbf{Applied Generative AI}, \textbf{Computational Finance}, \textbf{Open-Source Intelligence (OSINT)}, and \textbf{AI Systems for East African Markets}.
\end{itemize}

\section*{Referees}
\begin{itemize}
    \item \textbf{Mengi Bora Mtandama} \\
    Linguist and Chinese Translator \\
    Mbeya, Tanzania \\
    Phone: +255621906045

    \item \textbf{Andrew Malisa} \\
    Developer \\
    Azampay Tanzania \\
    Dar es Salaam \\
    Phone: +255767561957

    \item \textbf{Eva John Mwalyanga} \\
    TRA, Tanzania \\
    Phone: +255654678467
\end{itemize}

\section*{Declaration}
\noindent I declare that the information above is true to the best of my knowledge.

\vspace{10pt}
\noindent mwalyangashadrack@gmail.com \\
+255 625 534 921 \\
+255 774 174 921 \\
Dar es Salaam, Tanzania \\
\hlink{https://github.com/zuck30}{github.com/zuck30} \\
\hlink{https://antera.co.tz}{antera} \\
\hlink{https://zuck30.github.io/zuck30}{zuck30.github.io/zuck30}

\end{document}
